\chapter{Kết luận và Công việc tương lai\label{chapter:conclusions}}

Luận văn này so sánh C++ hiện đại và Rust trên một dịch vụ Quản lý Ngữ cảnh UE mô hình hoá theo
một thành phần gNB-CU-CP thực của Nokia, theo bốn chiều: độ dự đoán được của độ trễ đuôi, ngăn
chặn lỗi tại thời điểm biên dịch, chi phí mô hình an toàn, và chi phí áp dụng tăng dần. Xây dựng
trên một đặc tả cố định duy nhất và chạy qua cùng các đường mã, hai bản hiện thực cho phép tách
hiệu ứng ngôn ngữ khỏi khác biệt về thiết kế và kỹ năng.

Ba kết luận theo sau từ bằng chứng. Thứ nhất, về an toàn, khác biệt là khác biệt về \emph{loại}, và
nó có một cạnh xác định. Rust an toàn từ chối cả bảy lớp lỗi an toàn bộ nhớ và tranh chấp dữ liệu
được tiêm vào tại thời điểm biên dịch, và một trong bảy lớp đó là lớp mà không bộ dò lỗi C++ nào
báo cáo. Các mối nguy mà C++ chỉ có thể hy vọng bắt được trong một bài kiểm thử có gắn công cụ trở
thành không thể viết. Cạnh đó dừng ở đó.
Cả bảy lớp biên
đều biên dịch và hành xử sai trong Rust an toàn đúng như trong C++, và trên hai lớp trong số đó
công cụ C++ là bên báo cáo. Vì vậy bảo đảm được mô tả chính xác nhất như sau. Rust loại bỏ tranh
chấp dữ liệu, không phải điều kiện đua, và lỗi aliasing, không phải lỗi thứ tự khóa, rò rỉ, số học
quấn vòng hay ngăn chứa lỗi, và nó không loại bỏ lớp nào trong số đó ở vùng được đánh dấu
\texttt{unsafe}. Ở nơi các lỗi biên được bắt, chúng được
bắt tại thời điểm kiểm thử, trong hai trong ba trường hợp chỉ sau khi được viết lại theo các
primitive của bộ kiểm tra mô hình, đó là một chi phí thực và khác với chi phí của bộ dò lỗi.

Thứ hai, về di chuyển
chi phí chủ yếu do thiết kế API và đường cong học tập, không phải bởi bất kỳ thuế thời gian chạy
nào. Một lõi Rust không sửa đổi có thể được điều khiển từ C++ qua một đường nối mỏng, cô lập
\texttt{unsafe}, có chi phí trên một thao tác thực, đo so với cùng thao tác gọi bản địa, được giới
hạn ở $\FfiOpBoundPct\%$. Các chi phí mà kỹ sư build thường nêu trước không đổi thứ tự đó.
Phần lớn chi phí build thêm của Rust là biên dịch phụ thuộc chứ không phải công việc của chính
trình biên dịch, khoảng cách kích thước nhị phân là chính sách liên kết và đảo chiều khi chính sách
được khớp, và thành phần được so sánh không giải quyết mã bên thứ ba nào. Nơi C++ dẫn trước là vòng
lặp sửa-biên dịch, và nó chỉ dẫn trước với các sửa đổi nằm trong một đơn vị dịch.

Thứ ba, về hiệu năng bằng chứng chia đôi hai lần trước khi chia theo tài nguyên. Nó chia trước
theo \emph{cái gì đang được đo}, rồi theo \emph{liệu dịch vụ còn chịu tải được}. Làm đúng cả hai
phép chia là chỉnh sửa quan trọng nhất của nghiên cứu đối với chính nó, và mỗi phép được phát hiện
chứ không phải dự đoán trước. Một thao tác đầu-cuối trong thiết kế này là $\BudgetCoreShareMax\%$
công việc ngôn ngữ thực hiện, phần còn lại là bàn giao giữa luồng mà hệ điều hành thực hiện, nên
một kết quả null đầu-cuối không thể giới hạn hiệu ứng ngôn ngữ. Đo riêng thao tác thì hai ngôn ngữ
không tương đương. Rust nhanh hơn, theo $\BudgetCoreRatio\times$, vì
lựa chọn bố cục bảng băm của thư viện chuẩn chứ không vì điều gì
liên quan đến an toàn.

Phép chia thứ hai là theo chế độ, và nó quan trọng hơn. Dưới bão hoà hai ngôn ngữ tương đương ở
trung vị cho mọi mục đích triển khai quan tâm, và đó là kết quả thực về các triển khai dạng này.
Tương đương là thực tiễn chứ không phải thống kê, vì kiểm định ghép cặp vẫn phân giải
$\PFiftyDiffAbsMax\,\mu$s nghiêng về Rust dưới cả hai bộ cấp phát. Đuôi xa của chúng vẫn không phân
giải được trên máy ảo hoá, đó là thừa nhận về dụng cụ đo. Qua bão hoà chúng không tương đương. Bản
build C++ duy trì $\OverloadThrRatio\times$ thông lượng của bản build Rust, trong
\OverloadThrAgree{} lần lặp xen kẽ dưới cả hai bộ cấp phát, trong khi tiêu thụ
$\OverloadCpuPerOpRatio\times$ CPU trên mỗi thao tác hoàn thành, và các primitive cô lập không tái
tạo thứ tự đó ở mức ý nghĩa (Mục~\ref{sec:contention}). Về bộ nhớ bản build Rust mang thêm vài MiB
dấu chân cố định dưới cả hai bộ cấp phát, nhỏ hơn một bậc độ lớn so với chi phí của chính lựa chọn
bộ cấp phát.

Gộp lại, chúng ủng hộ một tuyên bố hẹp, hẹp hơn tuyên bố luận văn này dự kiến
đưa ra. Đối với các dịch vụ mặt điều khiển có trạng thái, đồng thời, có thể phân vùng dạng này, và
\emph{ở tải danh nghĩa}, lựa chọn giữa C++ hiện đại và Rust không phải quyết định hiệu năng. Đó là
quyết định về lớp lỗi nào tổ chức muốn làm không thể viết, và liệu tổ chức có trả chi phí kỹ thuật
và học tập để đạt được hay không. Dưới quá tải nó lại thành quyết định hiệu năng, và hiện nay ủng
hộ C++. Lý do nằm ở hạ tầng đồng thời chứ không phải bảo đảm của ngôn ngữ nào, và dịch vụ bắt buộc
phải suy giảm êm không thể coi chúng là phụ. Điều luận văn này bổ sung do đó là một dụng cụ đo chứ
không phải phán quyết: một đặc tả cố định, một cặp hiện thực đã kiểm thử tương đương, một khung đo
có thể tái lập vượt qua điểm gối bão hoà, một phân rã nói phần nào trong số đo đó ngôn ngữ phải
chịu trách nhiệm, và một danh mục lỗi ánh xạ cả phần bảo đảm bao phủ lẫn nơi nó dừng.

Một số hướng mở rộng công việc, mỗi hướng là biện pháp trực tiếp cho một hạn chế đã nêu ở trên.
Ưu tiên tức thời là chiến dịch bare-metal, mục tiêu giờ đã hẹp: điểm gối bão hoà và so sánh quá
tải đã được giải quyết, nên bare-metal sẽ mang lại đuôi xa dưới bão hoà. Bộ đếm phần cứng và hồ sơ
cấp phát thuộc cùng chuyến đo và là khoảng trống dụng cụ đo rõ nhất trong nghiên cứu. Chúng sẽ biến
kết quả bộ nhớ của Mục~\ref{sec:resources} từ quan sát thành giải thích, và xác nhận hoặc bác bỏ
lời giải thích về cục bộ bộ nhớ đệm mà Mục~\ref{sec:budget} đưa ra cho khoảng cách ở mức thao tác.
Ngoài ra:

\begin{itemize}
\item Phát lại các trace sản xuất đã ẩn danh hoá thay cho khối lượng công việc tổng hợp, để hỗn hợp
  thao tác được xác thực chứ không chỉ được biện minh bằng đặc tả.
\item Mở rộng so sánh sang thêm các phân hệ mặt điều khiển, hoặc sang đường L2/PHY nặng dữ liệu, nơi
  trạng thái dùng chung có thể thay đổi và áp lực cấp phát sẽ kiểm tra hai ngôn ngữ khác với thiết
  kế không chia sẻ.
\item Mở rộng danh mục lỗi dọc theo biên được ánh xạ ở đây. Lạm dụng \texttt{unsafe}, và
  các lớp biên dưới bộ kiểm tra mô hình toàn chương trình thay vì mô hình thu nhỏ thủ công, là các
  bước tiếp theo hiển nhiên. Một bộ kiểm tra bao phủ chương trình Rust cùng phụ thuộc nước ngoài
  của nó~\cite{pearce2026rustmc} sẽ chạm đường nối của
  Mục~\ref{sec:migrationresults}, mà nghiên cứu này đo chi phí nhưng không khám phá các xen kẽ của
  nó.
\item Đánh giá runtime async của Rust so với mô hình một luồng mỗi shard dùng ở đây, dưới các kiểm
  soát nghiên cứu này áp đặt. Điều đó sẽ tách sự phân kỳ đuôi mà Arvidsson và Kaller quy cho
  \texttt{async}/\texttt{await} không stack~\cite{arvidsson2025async} khỏi bản thân ngôn ngữ,
  và sẽ xác định bao nhiêu phần khoảng cách Olsson và Stenstr\"om đo giữa
  Rust bất đồng bộ và DPDK~\cite{olsson2023asyncpackets} thuộc runtime chứ không phải framework.
\item Đo chi phí cô lập trong tiến trình hai heap tại đường nối, vì lõi không \texttt{unsafe} cần
  được bảo vệ cũng như
  được xác minh~\cite{rivera2021galeed,kirth2022pkrusafe}.
\item Theo dõi mật độ lỗi và công sức review khi Rust được áp dụng tăng dần trong mã nguồn RAN sản
  xuất.
\end{itemize}
